\section{Giới thiệu tổng quan}\label{sec:introduction}

\subsection{Đặt vấn đề}
Dữ liệu giao thông đô thị phát sinh liên tục: mỗi chuyến xe bổ sung thông tin về điểm đón, quãng đường, chi phí và nhu cầu theo khu vực. Trong hệ thống giao thông thông minh (\emph{Intelligent Transportation System}, ITS), kết quả tổng hợp cần được cập nhật gần thời gian thực thay vì chỉ xuất hiện sau một đợt xử lý lô. Khi quy tắc chất lượng dữ liệu thay đổi, hệ thống cũng cần chuyển sang cách tính mới mà không làm gián đoạn truy vấn đang phục vụ.

Kho hồ dữ liệu (\emph{Data Lakehouse}) kết hợp tính linh hoạt của hồ dữ liệu với khả năng quản trị bảng và phân tích có cấu trúc của kho dữ liệu~\cite{armbrust2021lakehouse}. Khi bổ sung xử lý luồng (\emph{stream processing}), hệ thống có thể duy trì khung nhìn vật lý (\emph{materialized view}, MV), tức bảng kết quả được cập nhật theo dữ liệu mới thay vì tính lại toàn bộ mỗi lần truy vấn. Bài báo Ursa~\parencite{merli2025ursa} cho thấy nhu cầu thu hẹp khoảng cách giữa luồng sự kiện và lớp lưu trữ kho hồ dữ liệu. Từ động lực đó, đề tài tập trung vào một câu hỏi vận hành: trên cụm tài nguyên nhỏ, có thể phát lại dữ liệu, ghi đầu ra Iceberg, quan sát hệ thống và thay đổi cách tính MV công khai mà không ghi nhận lỗi truy vấn hay không?

\subsection{Lý do chọn đề tài}
Continux khảo sát câu hỏi trên bằng dữ liệu mở NYC TLC Yellow Taxi và các thành phần nguồn mở: K3s, Vector, Redpanda, RisingWave, MinIO, Apache Iceberg, Argo CD, VictoriaMetrics và Grafana. K3s giúp tái hiện các khái niệm cốt lõi của Kubernetes trên hạ tầng phòng thí nghiệm. Đề tài có liên hệ với Mục tiêu Phát triển Bền vững (\emph{Sustainable Development Goal}, SDG) số 8, 9 và 11 thông qua nhu cầu hạ tầng số hỗ trợ dịch vụ di chuyển hiệu quả, đổi mới công nghệ và đô thị bền vững~\parencite{sdg8,sdg9,sdg11}. Liên hệ này là động lực ứng dụng, không phải phép đo định lượng tác động xã hội.

\subsection{Mục tiêu nghiên cứu}
\begin{itemize}
\item Thiết kế luồng dữ liệu từ NYC TLC đến đầu ra Apache Iceberg trên MinIO.
\item Triển khai luồng trên cụm K3s ba máy chủ, quản lý trạng thái mong muốn bằng GitOps và quan sát bằng chỉ số vận hành.
\item Chạy bốn lượt phát lại độc lập với giới hạn phát cấu hình từ thấp đến cao.
\item Thực hiện chuyển giao Blue/Green (\emph{Blue/Green cutover}) ở lớp MV công khai và kiểm tra lỗi truy vấn trong lúc hoán đổi tên.
\item Chỉ kết luận từ dữ liệu có thể truy vết trong bộ bằng chứng đã lưu.
\end{itemize}

\subsection{Đối tượng và phạm vi}
Đối tượng xử lý là dữ liệu công khai NYC TLC Yellow Taxi tháng \texttt{2026-03} và bảng tra cứu Taxi Zone Lookup. Tệp Parquet được chuyển thành JSONL gồm \texttt{3.952.451} dòng; bảng tra cứu có \texttt{265} dòng. Mỗi lượt chạy tạo trạng thái nền sạch, phát lại dữ liệu trong cửa sổ đo cố định, chốt đầu ra, thực hiện chuyển giao MV và dọn trạng thái sinh ra.

Bốn cấu hình Vector là \texttt{smoke}, \texttt{benchmark-low}, \texttt{benchmark-medium} và \texttt{benchmark-high}, tương ứng với giới hạn phát cấu hình \texttt{2}, \texttt{1.000}, \texttt{5.000} và \texttt{10.000} sự kiện mỗi giây. Đây là giới hạn đặt tại nguồn phát, không phải thông lượng đầu cuối đã đo. Báo cáo sử dụng số chuyến đã xuất hiện trong MV, đường tăng theo thời gian, trạng thái cụm và đầu ra Iceberg để đánh giá hiện tượng quan sát được.

Phạm vi chuyển giao giới hạn ở tên truy vấn công khai \texttt{mv\_zone\_stats}. Báo cáo không khẳng định rằng đầu ra ghi Iceberg tự động chuyển quan hệ phụ thuộc sang cách tính mới sau phép hoán đổi tên MV.

\subsection{Câu hỏi nghiên cứu}
\begin{table*}[t!]
\centering
\caption{Câu hỏi nghiên cứu và bằng chứng trả lời}
\label{tab:rqs-vi}
\begin{tabularx}{\textwidth}{@{}p{0.08\textwidth}Y Y Y@{}}
\toprule
\textbf{ID} & \textbf{Câu hỏi} & \textbf{Bằng chứng} & \textbf{Kết luận trong phạm vi đo}\\
\midrule
RQ1 & MV công khai có tiếp tục truy vấn được trong lúc hoán đổi cách tính? & Bốn vòng lặp truy vấn, bốn số đo thời lượng và chỉ số lỗi. & Có; cả bốn lượt ghi nhận \texttt{0} lỗi truy vấn.\\
RQ2 & Luồng dữ liệu có tạo đầu ra kho hồ dữ liệu quan sát được? & Danh sách đối tượng Iceberg sau mỗi lượt phát lại. & Có; cả bốn lượt đều có Parquet và siêu dữ liệu Iceberg.\\
RQ3 & Quy trình có lặp lại được trên cụm tài nguyên giới hạn? & Trạng thái nút, PVC, khối lượng công việc, Argo CD, Vector và kết quả dọn dẹp. & Có; bốn lượt độc lập đều hoàn tất và trở về trạng thái sạch.\\
\bottomrule
\end{tabularx}
\end{table*}

\subsection{Đóng góp và cấu trúc báo cáo}
Đề tài đóng góp một hiện vật kỹ thuật có thể tái hiện, một quy trình thực nghiệm theo pha, bộ bằng chứng của bốn cấu hình phát và cách diễn giải thận trọng cho chuyển giao MV công khai. Chương 2 trình bày cơ sở lý thuyết. Chương 3 mô tả phương pháp và kiến trúc. Chương 4 trình bày kết quả. Chương 5 bàn luận giới hạn và kết luận. Phụ lục giữ SQL, cấu hình trích yếu và chỉ mục bằng chứng.
